% Incorporación aditiva. No modifica las fuentes ni el PDF de 211 páginas.
% Procedencia, dependencias y alcance en preservation_map.json.
\section{Representación grassmanniana de los lectores y formas con valores en la recta de acción}
\label{sec:compat-lectores-formas}

Las realizaciones positivas conservan información que puede volver a ser
utilizada por los lectores físicos. Esta conservación tiene una formulación
precisa: una representación marcada del registro permite reconstruir sus
coordenadas y transportar, por composición, cada función que las utiliza.
La geometría representa un resultado de la construcción; su elección de
coordenadas conserva la procedencia de ese resultado.

El dominio de partida está formado por los estados producidos por APP,
TRIT y TPK en las secciones precedentes. APP conserva sus hojas aditiva y
multiplicativa mediante residuo y cociente; TRIT selecciona el régimen
local con su orientación; la composición
\(U_t=\operatorname{Upd}_t\circ\operatorname{Tra}_t\circ
\operatorname{Sel}_t\) actualiza posición, registro y memoria.
La prolongación
\[
 w_6\longrightarrow w_{12}\longrightarrow w_{18}\longrightarrow
 w_{24}\longrightarrow w_{30}\longrightarrow R_{36}\longrightarrow G_9
\]
conserva el prefijo y distingue retorno de fase de retorno del estado.
Las realizaciones espectral, solenoidal, gauge, cohomológica y determinantal
permanecen como operaciones correlativas de la estructura discreta
conjunta del continuo. La extracción del registro y las representaciones
regionales se realizan en ese mismo dominio. En lo que sigue, todas ellas
son salidas HMT previamente construidas.

\subsection{La colección marcada y su inversa}

Sea \(L\geq2\), y sea \(K=(K_1,\ldots,K_L)\) un registro positivo
con carga \(Q=\sum_jK_j\). En la aplicación dodecafásica,
\(L=12\) y \(Q=6263\). Definimos
\[
 d_j=K_j/Q,\qquad t_0=0,\qquad
 t_j=\sum_{a=1}^{j}d_a,\qquad
 B(K)=\begin{pmatrix}1&\cdots&1\\t_0&\cdots&t_L\end{pmatrix}.
\]
La colección \(\mathcal G_L^{\rm aff}\) es la imagen de estos planos
de filas en \(\operatorname{Gr}_{>0}(2,L+1)\), conservando las
columnas ordenadas \(0,\ldots,L\), las marcas extremas y la carga
\(Q\). La normalización afín fija \(t_0=0,t_L=1\).
Escribimos \(\Delta_{ij}\) para el menor de las columnas \(i,j\)
de cualquier representante del plano.

\begin{proposition}[Inversa de la realización marcada]
\label{prop:compat-inversa-marcada}
La aplicación \(K\mapsto([B(K)],Q)\) es inyectiva. Su inversa
sobre la imagen está dada por
\begin{equation}
 K_j=Q\frac{\Delta_{j-1,j}}{\Delta_{0L}},
 \qquad
 t_j=\frac{\Delta_{0j}}{\Delta_{0L}},
 \quad 1\leq j\leq L.
 \label{eq:compat-inversa-plucker}
\end{equation}
Los cocientes son invariantes bajo cualquier cambio invertible de base
de las dos filas.
\end{proposition}
\begin{proof}
En el representante declarado,
\(\Delta_{ij}=t_j-t_i\), y \(\Delta_{0L}=1\).
Las fórmulas recuperan las diferencias consecutivas y sus sumas.
Si se sustituye \(B\) por \(AB\), con \(A\) invertible,
todos los menores se multiplican por \(\det A\); sus cocientes
permanecen iguales. La carga recupera la escala anterior a la
normalización. Por tanto, dos registros con iguales datos marcados
tienen las mismas coordenadas.
\end{proof}

La condición de pertenencia a \(\mathcal G_L^{\rm aff}\) importa.
Los cocientes consecutivos de un plano positivo arbitrario pueden
incumplir la identidad \(\sum_j\Delta_{j-1,j}=\Delta_{0L}\).
En la colección construida esa identidad es telescópica. La inversa actúa
en la imagen marcada que acaba de definirse.

\subsection{Álgebra de lectores y cambios de coordenatización}

Sea \(\mathcal D\) la imagen de una representación del estado enriquecido
de la forma \(x\mapsto(K(x),\xi(x))\). El símbolo \(\xi\)
reúne exactamente las restantes coordenadas utilizadas por el lector:
representaciones regionales, origen y orientación, ruta, firma, condiciones
de frontera y secciones dimensionales, según el objeto considerado.
No se supone que estas coordenadas sean independientes. Definimos
\[
 \mathcal F:\mathcal D\longrightarrow\widetilde{\mathcal G},
 \qquad (K,\xi)\longmapsto([B(K)],Q,\xi),
 \qquad \widetilde{\mathcal G}:=\mathcal F(\mathcal D).
\]
La proposición anterior proporciona la inversa \(\mathcal I\) de
\(\mathcal F\) en su imagen.

\begin{theorem}[Representación fiel del álgebra de lectores retenidos]
\label{thm:compat-algebra-lectores}
Para cada lector \(\mathcal R:\mathcal D\to V\), con codominio
declarado \(V\), existe un único lector geométrico
\(\mathcal R^{\rm geom}:\widetilde{\mathcal G}\to V\) tal que
\[
 \mathcal R^{\rm geom}=\mathcal R\circ\mathcal I,
 \qquad \mathcal R^{\rm geom}\circ\mathcal F=\mathcal R.
\]
Para lectores escalares, esta correspondencia preserva sumas, productos,
unidad y conjugación. Es un isomorfismo del álgebra de funciones sobre
\(\mathcal D\) con el álgebra correspondiente sobre
\(\widetilde{\mathcal G}\).
\end{theorem}
\begin{proof}
La composición está definida porque \(\mathcal I\mathcal F\) es
la identidad. Si dos lectores geométricos tienen la misma composición
con \(\mathcal F\), coinciden en su imagen, que es todo
\(\widetilde{\mathcal G}\). Las operaciones algebraicas se verifican
punto a punto; por ejemplo,
\((\mathcal R_1\mathcal R_2)\circ\mathcal I=
(\mathcal R_1\circ\mathcal I)(\mathcal R_2\circ\mathcal I)\).
La composición inversa es la restricción por \(\mathcal F\).
\end{proof}

La fidelidad se refiere a la representación \((K,\xi)\). Para un
observable \(\mathcal O\) del estado enriquecido completo, el criterio
de descenso es que \(\mathcal O(x)=\mathcal O(x')\) siempre que
\((K(x),\xi(x))=(K(x'),\xi(x'))\). La necesidad se obtiene por
composición; la suficiencia permite definir el valor en una fibra
mediante cualquiera de sus representantes. Este criterio conserva
explícitamente qué información debe acompañar al registro.

Sea ahora \(\mathcal T\) una triangulación del polígono etiquetado
que determina una coordenatización de Plücker positiva. Se incluyen sus coordenadas
de frontera y una normalización homogénea, por ejemplo
\(\Delta_{0L}=1\). Si se intercambia la diagonal \(ac\) por
\(bd\), con \(a<b<c<d\), el cambio de coordenatización es
\begin{equation}
 \Delta_{bd}=
 \frac{\Delta_{ab}\Delta_{cd}+\Delta_{ad}\Delta_{bc}}{\Delta_{ac}}.
 \label{eq:compat-cambio-carta}
\end{equation}
El denominador es positivo y la relación de Plücker prueba que ambas
coordenatizaciones evalúan el mismo punto.

\begin{corollary}[Compatibilidad de los lectores con las mutaciones]
\label{cor:compat-mutaciones}
Sean \(r_{\mathcal T}\) y \(r_{\mathcal T'}\) las expresiones
de \(\mathcal R^{\rm geom}\) en dos coordenatizaciones que conservan las marcas,
y \(\mu_{\mathcal T\mathcal T'}\) su cambio positivo. Entonces
\[
 r_{\mathcal T'}\circ\mu_{\mathcal T\mathcal T'}
 =r_{\mathcal T}
\]
en el dominio común. Una sucesión cerrada de cambios de coordenatización que
conserva el punto marcado conserva todos estos lectores.
\end{corollary}
\begin{proof}
Las dos reconstrucciones mediante
\eqref{eq:compat-inversa-plucker} recuperan el mismo \(K\).
Las coordenadas \(\xi\) se transportan con sus marcas. Aplicar
\(\mathcal R\) prueba la igualdad; iterarla prueba la afirmación
sobre una sucesión cerrada.
\end{proof}

Los lectores pueden expresarse en coordenatizaciones distintas y seguir siendo
funciones del mismo objeto. Una permutación de las etiquetas actúa
adicionalmente sobre la marca posicional; una transformación que cambia
el punto actúa sobre el argumento del lector. Ambas operaciones tienen
su ley de transporte específica.

La acción diagonal por columnas proporciona un control explícito.
Para \(\lambda_i>0\),
\(\Delta_{ij}\mapsto\lambda_i\lambda_j\Delta_{ij}\).
La razón doble \(\Delta_{ab}\Delta_{cd}/
(\Delta_{ad}\Delta_{bc})\) se conserva, mientras
\(\Delta_{j-1,j}/\Delta_{0L}\) recibe el factor
\(\lambda_{j-1}\lambda_j/(\lambda_0\lambda_L)\).
Con \(t=(0,1/4,1/2,1)\) y \(\lambda=(1,2,1,1)\), las
separaciones reconstruidas serían \((1/2,1/2,1/2)\), cuya suma
es \(3/2\). Las razones dobles solas conservan un cociente diferente
del requerido por la inversa marcada.

\subsection{Composición explícita con acción y respuesta constitutiva}

En la coordenatización dodecafásica, escribimos \(b=1000\) y
\begin{equation}
 \kappa^{\rm geom}=
 \frac{Q}{b^{12}-1}\sum_{j=1}^{12}
 b^{12-j}\frac{\Delta_{j-1,j}}{\Delta_{0,12}},
 \qquad
 \alpha^{\rm geom}=\pi_{\rm HMT}+e_{\rm HMT}
 -\varphi_{\rm HMT}-4-\kappa^{\rm geom}.
 \label{eq:compat-alpha-geometrica}
\end{equation}
Son las expresiones geométricas del lector de clausura
\eqref{eq:mass-k-alpha}. Las tres representaciones regionales y \(K\)
comparten generación coinductiva; la expresión utiliza su orden
constructivo interno. El lector de acción
\eqref{eq:mass-k-action}, evaluado en \(\alpha^{\rm geom}\),
determina \(\eta_{\rm ret}^{\rm geom}\), y por tanto
\[
 \hbar_{\rm ret}^{\rm geom}
 =\eta_{\rm ret}^{\rm geom}10^{-34}\mathcal U_S
 =\hbar_{\rm ret}.
\]
La igualdad tiene lugar en la recta de acción \(\mathcal L_S\)
con su sección \(\mathcal U_S\); conserva la década procedente
del lector determinantal \(\varphi_{\rm HMT}\alpha^{16}\).
La composición no modifica la sección ni la aplica una segunda vez.

En la coordenatización angular levantada que conserva el origen y la rama real
de la representación, el par angular posterior determina
\[
 x=\frac{\pi_{\rm HMT}}{180}(1000\alpha),\qquad
 y=\frac{\pi_{\rm HMT}}{90}(\eta_{\rm ret}+\alpha),\qquad
 q_\pm=e^{-x\mp y}.
\]
La rama forma parte de las marcas retenidas: la reducción de un ángulo
módulo \(360\) por sí sola conserva otra información. Las igualdades
que siguen emplean el levantamiento declarado, en el cual la primera
coordenada es \(1000\alpha\).
En la cámara \(x>|y|\), con proyectores de hoja fijados
\(P_\pm\), sea \(T=q_+P_++q_-P_-\). La respuesta del corpus
compara las incidencias \(90\) y \(120\) mediante
\[
 \mathcal V=(I-T^{90})(I-T^{120})^{-1}
 =r_+P_++r_-P_-,\qquad
 r_\pm=f(q_\pm^{30}),\quad
 f(s)=\frac{1+s+s^2}{1+s+s^2+s^3}.
\]
La regla satisface
\(\mathcal V\Sigma_4(T^{30})=\Sigma_3(T^{30})\), donde
\(\Sigma_j(S)=\sum_{a=0}^{j-1}S^a\). Su unicidad se deduce
de la invertibilidad de \(\Sigma_4(T^{30})\).
Las cuatro lecturas normalizadas son
\[
 \widehat\varepsilon=r_+^2,\quad
 \widehat\mu=r_-^2,\quad
 \widehat Z=r_-/r_+,\quad
 \widehat c=(r_+r_-)^{-1}.
\]
En las rectas dimensionales de acción, carga y velocidad se realizan como
\[
 Z=\widehat Z\,u_Su_Q^{-2},\quad c=\widehat c\,u_C,
 \qquad
 \mu=\widehat\mu\,u_Su_C^{-1}u_Q^{-2},\quad
 \varepsilon=\widehat\varepsilon\,u_S^{-1}u_C^{-1}u_Q^2.
\]
Las identidades \(\mu\varepsilon=c^{-2}\) y
\(\mu/\varepsilon=Z^2\) se obtienen por producto y cociente.
El corolario de cambios de coordenatización conserva esta respuesta completa,
con las marcas de hoja y las secciones dimensionales incluidas en
\(\xi\). El mismo argumento se aplica a un lector de masa con su
ruta, firma y operador de retorno retenidos, en el dominio de la
sección~\ref{sec:k-mass-bridge}.

\begin{theorem}[Recuperación del registro desde la respuesta etiquetada]
\label{thm:compat-recuperacion-constitutiva}
Fijadas las representaciones regionales, la orientación de las hojas, la
base \(b=1000\), el origen, la longitud doce y el levantamiento angular
declarado, la aplicación
\(K\mapsto(\widehat\varepsilon,\widehat\mu)\) es inyectiva
sobre los registros enteros de su imagen física que satisfacen
\(0\leq K_j<b\) y la cámara contractiva indicada. Su inversa exacta
es la composición de las siguientes operaciones:
\begin{align*}
 r_+&=\sqrt{\widehat\varepsilon},&
 r_-&=\sqrt{\widehat\mu},\\
 r_\pm s_\pm^3+(r_\pm-1)(1+s_\pm+s_\pm^2)&=0,
 &0<s_\pm<1,\\
 x&=-\frac1{60}\log(s_+s_-),&
 \alpha&=\frac{180x}{1000\pi_{\rm HMT}},\\
 \kappa&=\pi_{\rm HMT}+e_{\rm HMT}-\varphi_{\rm HMT}-4-\alpha,
 &N&=(b^{12}-1)\kappa,\\
 K_j&=\left\lfloor N/b^{12-j}\right\rfloor\bmod b.
\end{align*}
\end{theorem}
\begin{proof}
Las raíces positivas recuperan los coeficientes etiquetados.
La derivada
\[
 f'(s)=-\frac{s^2(3+2s+s^2)}{(1+s+s^2+s^3)^2}<0
 \qquad(0<s<1)
\]
y sus límites \(1\) y \(3/4\) prueban que la ecuación cúbica
posee exactamente una raíz en ese intervalo. Puesto que
\(s_\pm=e^{-30x\mp30y}\), su producto es \(e^{-60x}\).
Se recuperan \(x\), \(\alpha\) y \(\kappa\).
En la imagen considerada, \(N=\sum_jK_jb^{12-j}\) es un entero
entre cero y \(b^{12}-1\). La división euclídea sucesiva recupera
de manera única sus doce cifras, incluidos los ceros iniciales.
\end{proof}

La precisión necesaria para una recuperación numérica también puede
expresarse. Con representaciones regionales exactas, un error
\(|\widetilde\alpha-\alpha|<1/[2(b^{12}-1)]\) implica
\(|\widetilde N-N|<1/2\); el redondeo al entero más próximo
recupera \(N\). Si las representaciones regionales tienen errores, se
aplica la misma condición a la suma de errores de
\(\pi+e-\varphi-\alpha\). La inyectividad exacta y la resolución
metrológica son así propiedades diferenciadas.

El dominio entero es esencial. En una prolongación analítica de los
registros positivos, la perturbación de tres posiciones consecutivas
proporcional a \((1,-(b+1),b)\) conserva tanto la suma como la
lectura posicional, pues
\(1-(b+1)+b=0\) y \(b^2-b(b+1)+b=0\).
Para perturbaciones suficientemente pequeñas conserva la positividad
y cambia el punto grassmanniano. La inyectividad precedente utiliza
la condición de doce cifras enteras; la colección analítica continua
posee otras fibras.

\subsection{Refinamiento y lectores heredados}

Si cada separación \(d_j\) se subdivide en separaciones positivas
\(d_{j,a}\) con \(\sum_ad_{j,a}=d_j\), los extremos originales
conservan su posición. Sea \(\rho\) la operación de sumar los sucesores
de cada intervalo y \(H\) la inclusión de sus extremos en el nuevo
conjunto ordenado. Entonces
\[
 \mathcal I_{\rm heredada}(B(\widetilde d),Q)
 =\bigl(Q\sum_ad_{j,a}\bigr)_j=K,
 \qquad
 \mathcal R_{\rm heredada}=\mathcal R(K,\xi).
\]
La igualdad prueba la conmutatividad entre reconstrucción heredada y
subdivisión. En una sucesión de subdivisiones, las sumas se asocian;
el resultado persiste en cada composición finita. Esto conserva el
lector de longitud doce y sus marcas. Un lector definido sobre los
nuevos bloques tendrá su propio dominio y su ley de comparación.

\subsection{Formas canónicas con valores en una recta dimensional}

Sea \(P\) uno de los politopos de rango uno construidos en el
manuscrito y \(P=\bigcup_\sigma P_\sigma\) una subdivisión
orientada del mismo soporte, con multiplicidad uno y caras compatibles.
Escribimos \(\Omega_\sigma\) y \(\Omega_P\) para sus formas
canónicas. Sus identidades son
\(\Omega_P=\sum_\sigma\Omega_\sigma\) y
\(\operatorname{Res}_F\Omega_P=\Omega_F\), con la convención
de orientación fijada en el desarrollo amplituhedral.
La recta de acción \(\mathcal L_S\) se considera constante sobre
este politopo; su sección procede del estado HMT fijado.

\begin{theorem}[Transporte de la forma por una sección de acción]
\label{thm:compat-forma-accion}
Para \(a\in\mathcal L_S\), constante respecto de las coordenadas
del politopo,
\[
 \boldsymbol\Omega_P:=a\otimes\Omega_P
 =\sum_\sigma a\otimes\Omega_\sigma,
 \qquad
 \operatorname{Res}_F\boldsymbol\Omega_P=a\otimes\Omega_F.
\]
Las identidades se conservan bajo subdivisiones sucesivas y residuos
iterados sobre banderas de caras. Se aplican, en particular, a
\(a=\hbar_{\rm ret}\).
\end{theorem}
\begin{proof}
El producto tensorial con una recta es lineal. El residuo actúa sobre
el factor diferencial y conmuta con el coeficiente constante. La
compatibilidad de subdivisiones y de residuos iterados se transporta
por las mismas operaciones lineales. En un cambio de base dimensional
\(u'_S=\lambda\,u_S\), con \(\lambda>0\), el coeficiente numérico
cambia por \(\lambda^{-1}\);
el tensor y sus residuos permanecen iguales.
\end{proof}

\begin{theorem}[Condición de pegado de los coeficientes de frontera]
\label{thm:compat-pegado}
Sean \(a_\sigma\) secciones de la misma recta, regulares en un
entorno de las caras y meromorfas en una realización ambiental común.
Supóngase que sus productos con las formas canónicas tienen, a lo
largo de las caras consideradas, polos a lo sumo simples. En el
interior relativo de una cara común \(F\) de codimensión uno de dos celdas
\(\sigma,\tau\), con orientaciones opuestas,
\begin{equation}
 \operatorname{Res}_F\left(a_\sigma\Omega_\sigma+a_\tau\Omega_\tau\right)
 =\left(a_\sigma|_F-a_\tau|_F\right)\Omega_F.
 \label{eq:compat-defecto-frontera}
\end{equation}
Por tanto, el polo de esta pareja sobre el divisor de \(F\) desaparece
exactamente cuando las dos trazas coinciden. Si los coeficientes son
constantes y el grafo de adyacencia de las celdas es conexo, la cancelación
de cada pareja de celdas a través de su cara común equivale a la existencia
de un único coeficiente común \(a\).
\end{theorem}
\begin{proof}
Con coordenada local transversal \(z\), una forma logarítmica se
escribe \(dz/z\wedge\omega+\eta\), donde \(\eta\) es regular
respecto de ese divisor. La regularidad de \(a_\sigma\) da
\(\operatorname{Res}_F(a_\sigma\Omega_\sigma)
=a_\sigma|_F\operatorname{Res}_F\Omega_\sigma\).
Las dos orientaciones producen \(\Omega_F\) y \(-\Omega_F\).
La forma \(\Omega_F\) es no nula en el interior relativo de la
cara; la igualdad de residuos equivale a la igualdad de las trazas.
Para un polo simple, la anulación del residuo elimina el polo.
En el caso constante, la igualdad se propaga a lo largo de cada
arista del grafo conexo de adyacencia.
\end{proof}

Para la suma global deben contarse todos los términos con un polo en el
mismo divisor ambiental. Si \(H\) es el hiperplano que contiene \(F\),
y \(I_H\) reúne esos términos, la identidad general es
\[
 \operatorname{Res}_{H}\left(\sum_\nu a_\nu\Omega_\nu\right)
 =\sum_{\nu\in I_H}a_\nu|_H\operatorname{Res}_H\Omega_\nu,
\]
siempre que los coeficientes sean regulares sobre \(H\). Una faceta
coplanar, aunque no sea adyacente a la cara real \(F\), participa en esta
suma. La fórmula de dos términos representa el residuo global cuando
los restantes términos son regulares sobre ese divisor. La cancelación
por parejas de todas las facetas interiores elimina sus contribuciones;
las facetas exteriores conservan los residuos de frontera correspondientes.

En una cara de codimensión superior, se aplica esta fórmula sucesivamente
con las orientaciones inducidas. Los signos son los de la incidencia
celular y satisfacen la cancelación de borde de borde. La colección de
lectores locales define un pegado compatible por parejas cuando sus
trazas coinciden después del transporte a la misma recta de coeficientes.
El residuo global se obtiene sumando las contribuciones de cada divisor.
Este enunciado da un criterio efectivo para distribuir una lectura
física sobre coordenatizaciones o celdas sin introducir discontinuidades espurias.

\subsection{Cancelación de polos y normalización global}

La cancelación interior determina una condición local. La identificación
con una forma canónica global requiere asimismo el divisor de polos y
los residuos exteriores normalizados. En espacio proyectivo, dos
formas racionales de grado máximo con los mismos polos logarítmicos
y los mismos residuos difieren por una forma holomorfa de grado máximo;
\(H^0(\mathbb P^m,\Omega^m)=0\) demuestra su igualdad.
Ésta es la condición global utilizada en la prueba del politopo
\cite{positive}.

Un ejemplo intervalar muestra por qué los coeficientes variables
requieren dicha condición. Para \(0<c<1\), sean
\[
 \Omega_{0c}=\frac{c\,dx}{x(c-x)},\qquad
 \Omega_{c1}=\frac{(1-c)\,dx}{(x-c)(1-x)},\qquad
 \Omega_{01}=\frac{dx}{x(1-x)}.
\]
Con \(a_1(x)=1+x(c-x)\) y \(a_2(x)=1\), sus trazas coinciden
en todos los extremos pertinentes, y sin embargo
\[
 a_1\Omega_{0c}+a_2\Omega_{c1}=\Omega_{01}+c\,dx.
\]
El término adicional es regular en la coordenatización afín y tiene un polo en
el infinito proyectivo. Los residuos de los extremos finitos y la
cancelación en \(c\) se conservan, mientras la forma resultante
tiene otro comportamiento global. Multiplicar el ejemplo por una
sección de \(\mathcal L_S\) produce el mismo control dimensional.

La consecuencia de conjunto es precisa: la representación grassmanniana
marcada transporta los lectores retenidos; las mutaciones conservan sus
valores cuando sólo cambian la coordenatización; el refinamiento conserva sus
evaluaciones heredadas; y las formas con valores en la recta de acción
se ensamblan mediante igualdad de trazas y control de polos. Cada
igualdad expresa la operación que conserva el objeto correspondiente.
